%% ============================================================================
%% revised_abstract_v2.tex
%%
%% Round-2 follow-up: trimmed to <=220 words (target: 200-215).
%% Same calibrated novelty as v1, but with tightened phrasing in P1+P2 and a
%% compressed final paragraph that drops the "no claim of new wild character
%% variety theorems" sentence (moved to Section 1.5 of the paper).
%%
%% Word count (non-math, EQ display block counted as 1): 209
%% ============================================================================
\begin{abstract}
For a scalar local linear operator $L$ (differential, difference, or
$q$-difference) with one irregular singularity at $x = \infty$ whose
Newton polygon has a unique dominant edge of slope $-p/q$ with $r{+}1$
lattice points, the WKB ansatz on that edge produces a characteristic
polynomial $\chi \in \mathbb{C}[c]$. Our central algebraic identity is
\[
   \chi(c) \;=\; \chi_w(c^q),
\]
where $\chi_w \in \mathbb{C}[w]$ is the classical edge-indicial polynomial
of degree $r$ in $w = c^q$. The exponent constellation
$\Cl(L) \subset \mathbb{C}$ is therefore a $\mu_q$-equivariant union of
$r$ fibres of $c \mapsto c^q$ over $\Roots(\chi_w)$.

This factorisation organises a finite stratified atlas: a computational
sweep yields $23$ single-edge classes for
$-p/q \in \{1/2,\,1/3,\,1/4,\,2/3,\,3/2\}$ and $r \le 3$, refining to
$39$ once multi-edge, nested-ramification, and prime-$q$ arithmetic
strata are tracked. We prove a per-edge factorisation along the slope
filtration, an iterated $\mu_{q'}$-pullback law for nested inner
polynomials, a block-diagonal pullback for systems, and a prime-$q$
Kummer dichotomy separating the algebraically independent stratum from
the cyclotomic stratum.

The analytic input attaching anti-Stokes directions is classical
(Levelt--Turrittin, Sibuya, Wasow); on the $q$-fold cover the rays of a
generic single-edge operator satisfy
$\arg \Dij + p\phi \equiv \pi/2 \pmod{\pi}$. One explicit consequence:
the slope numerator $p$, invisible to $\Cl(L)$, is visible to the
cover-side ray pattern. The analytic statements
(Theorems~\ref{thm:AT1}, \ref{thm:AT3}, \ref{thm:AT5}, \ref{thm:AT7})
are $\mu_q/\mu_{q'}$-equivariant repackagings of classical
formal-classification theory.
\end{abstract}
